\chapter{Measure-Theoretic Lemmas and Integrability Witnesses}
\label{app:measure-lemmas}

This appendix proves the measure-theoretic statements of \Cref{ch:probability,ch:elbo} and separates the
hypotheses of \cref{hyp:elbo-evidence-domain} from factorwise log-integrability
(\cref{hyp:elbo-factorwise-integrability}) by witnesses.

\section{Kernel measurability and proofs for the probability chapter}

\theoremheading{Jointly measurable Radon--Nikodym version for a dominated kernel}{thm:prob-kernel-rn-measurable-version}
Let $(\mathsf E,\mathscr E)$ be a measurable space, and let $K$ and $L$ be probability kernels from $\mathsf E$
to $(\mathsf F,\mathscr F)$ with $K(x,\cdot)\ll L(x,\cdot)$ for every $x$. Then there is a finite
$\mathscr E\otimes\mathscr F$-measurable $r:\mathsf E\times\mathsf F\to[0,\infty)$ with
$K(x,B)=\int_Br(x,y)L(x,dy)$ for every $x\in\mathsf E$ and $B\in\mathscr F$. Selecting a derivative separately
for each $x$ does not give joint measurability; the theorem supplies one simultaneous version.
\status{ESTABLISHED}

\noindent\emph{Proof.}
Standard Borelness gives finite refining measurable partitions $(\mathcal P_n)_{n\geq1}$ whose generated
sigma-algebras increase to $\mathscr F$. With $0/0:=0$ put
\begin{equation}
r_n(x,y)=\sum_{A\in\mathcal P_n}\frac{K(x,A)}{L(x,A)}\mathbf 1_A(y),
\label{eq:prob-kernel-rn-partitions}
\end{equation}
and let $r=\limsup_nr_n$ where this is finite and $r=0$ elsewhere; like every $r_n$, $r$ is jointly
measurable. For fixed $x$, $r_n(x,\cdot)$ is the conditional expectation of $h_x=dK(x,\cdot)/dL(x,\cdot)$ given
$\sigma(\mathcal P_n)$ under $L(x,\cdot)$, so martingale convergence gives $r_n(x,\cdot)\to h_x<\infty$
$L(x,\cdot)$-almost everywhere, and $r(x,\cdot)$ is a derivative for every $x$. \hfill$\square$

\corollaryheading{Joint densities against one reference}{cor:prob-common-reference-joint-density}
Let $\nu$ be a sigma-finite measure on $(\mathsf F,\mathscr F)$ and $K$ a probability kernel from
$(\mathsf E,\mathscr E)$ to $\mathsf F$ with $K(x,\cdot)\ll\nu$ for every $x$. Then $K$ has a finite jointly
measurable density against $\nu$, and every finite or countable collection of kernels dominated by the same
$\nu$ has jointly measurable densities against that common reference; in particular this applies to the
generative and recognition kernels \eqref{eq:prob-generative-signature}--\eqref{eq:prob-recognition-signature}
when they are dominated by the product references. \status{ESTABLISHED}

\noindent\emph{Proof.}
Sigma-finiteness supplies measurable $w$ with $0<w<\infty$ everywhere and $\int w(y)\nu(dy)=1$, so the constant
kernel $L(x,dy)=w(y)\nu(dy)$ is equivalent to $\nu$. \Cref{thm:prob-kernel-rn-measurable-version} gives a finite
jointly measurable $r=dK/dL$, and $k=rw$ is the density. Apply this to each member of the collection.
\hfill$\square$

\propositionheading{Kernel integration preserves measurability}{prop:prob-kernel-integration-measurability}
Let $\kappa$ be an s-finite kernel from $(\mathsf E,\mathscr E)$ to $(\mathsf F,\mathscr F)$, that is, a
countable sum of finite kernels; probability kernels and fixed sigma-finite measures are instances. For
every $\mathscr E\otimes\mathscr F$-measurable $f:\mathsf E\times\mathsf F\to[0,\infty]$,
\label{prop:prob-sigmafinite-integration-measurability}
\begin{equation}
x\longmapsto\int_{\mathsf F}f(x,y)\kappa(x,dy)
\label{eq:prob-kernel-integration}
\end{equation}
is $\mathscr E$-measurable. For an extended signed integrand the same holds on the measurable set where the
positive-part and negative-part integrals are not both infinite, and no expression $+\infty-\infty$ is formed.
Consequently the evidence \eqref{eq:prob-evidence} and the $\mathbb Q_{V,o,X}$-expectations of \Cref{ch:elbo}
are measurable in $(o,X)$ whenever their integrands are jointly measurable. \status{ESTABLISHED}

\noindent\emph{Proof.}
For finite $\kappa$, the class of $C\in\mathscr E\otimes\mathscr F$ for which $x\mapsto\kappa(x,C_x)$ is
measurable, with $C_x=\{y:(x,y)\in C\}$, contains the rectangles, since
$\kappa(x,(A\times B)_x)=\mathbf 1_A(x)\kappa(x,B)$, and is closed under proper differences and increasing
unions because each $\kappa(x,\cdot)$ is finite. Rectangles form a $\pi$-system generating
$\mathscr E\otimes\mathscr F$, so the $\pi$--$\lambda$ theorem gives every $C$. Linearity and monotone
convergence pass to nonnegative measurable $f$, and countable sums to s-finite $\kappa$; a sigma-finite $\nu$
with a measurable partition into sets $F_m$ of finite measure is the constant s-finite kernel
$\sum_m\nu(\cdot\cap F_m)$. A signed integrand is split into $f^+$ and $f^-$, whose difference is measurable
where they are not both infinite. \hfill$\square$

\corollaryheading{Measurability of kernel relative entropy}{cor:prob-kernel-kl-measurability}
Let $\kappa$ and $\eta$ be probability kernels from $(\mathsf E,\mathscr E)$ to $(\mathsf F,\mathscr F)$. Then
$x\mapsto\KL(\kappa(x,\cdot)\Vert\eta(x,\cdot))$ is $\mathscr E$-measurable as a $[0,+\infty]$-valued map, and
for finite refining partitions $(\mathcal P_n)$ generating $\mathscr F$,
\begin{equation}
\KL(\kappa(x,\cdot)\Vert\eta(x,\cdot))
=\sup_{n\geq1}\sum_{A\in\mathcal P_n}\kappa(x,A)\log\frac{\kappa(x,A)}{\eta(x,A)},
\label{eq:prob-kernel-kl-partitions}
\end{equation}
with $0\log(0/b)=0$ and $a\log(a/0)=+\infty$ for $a>0$. \status{ESTABLISHED}

\noindent\emph{Proof.}
Put $L=(\kappa+\eta)/2$. \Cref{thm:prob-kernel-rn-measurable-version} applied to $(\kappa,L)$ and $(\eta,L)$
gives finite jointly measurable densities $a$ and $b$. Fix $x$, write $L_x=L(x,\cdot)$, and let $a_n,b_n$ be the
conditional expectations of $a(x,\cdot),b(x,\cdot)$ under $L_x$ given $\sigma(\mathcal P_n)$, whose cell values
are the cellwise ratios in \eqref{eq:prob-kernel-kl-partitions}. The perspective
$\phi_{\mathrm{KL}}(u,v)=u\log(u/v)-u+v$, with $\phi_{\mathrm{KL}}(0,v)=v$, $\phi_{\mathrm{KL}}(u,0)=+\infty$ for
$u>0$, and $\phi_{\mathrm{KL}}(0,0)=0$, is nonnegative, convex, and lower semicontinuous. Since both kernels have
unit mass, the $n$-th partition sum is $D_n(x)=\int\phi_{\mathrm{KL}}(a_n,b_n)dL_x$ and the divergence is
$\int\phi_{\mathrm{KL}}(a,b)dL_x$. Conditional Jensen makes $D_n(x)$ nondecreasing in $n$ and at most the
divergence, and martingale convergence $(a_n,b_n)\to(a,b)$ $L_x$-almost everywhere, with Fatou's lemma and
lower semicontinuity, gives $\int\phi_{\mathrm{KL}}(a,b)dL_x\leq\liminf_nD_n(x)$. Each $D_n$ is measurable in
$x$, so the countable supremum is measurable. \hfill$\square$

\noindent\emph{Proof of \Cref{thm:prob-mixed-family-common-domination}.}
If $A_*$ is countable, $\nu_{\mathrm{mix}}$ is sigma-finite, and a $\nu_{\mathrm{mix}}$-null set is Lebesgue null
and contains no point of $A_*$, so every term of every $\mu_t$ vanishes on it. Conversely, let a sigma-finite
$\xi$ dominate every $\mu_t$. Each $a\in A_*$ has $\mu_t(\{a\})>0$ for some $t$, whence $\xi(\{a\})>0$, and a
sigma-finite measure has at most countably many points of positive mass, as in the proof of
\Cref{sec:exp-dominated-nondominated}. The moving-atom family has $A_*=\R$. \hfill$\square$

\remarkheading{A regular-conditional version is not determined at a point}{rem:prob-version-witness}
On $\R^2$ with coordinates $(o,y)$, let $\phi$ be the standard normal density and
$\psi\neq\phi$ another probability density. Then $p(o,y)=\phi(o)\phi(y)$ and $\widetilde p$, equal to $p$ for
$o\neq0$ and to $\phi(0)\psi(y)$ for $o=0$, are versions of the same joint density; both evidence
representatives equal $\phi(0)$ at $o=0$, but their density-ratio conditionals there are $\phi(y)dy$ and
$\psi(y)dy$. The joint measure therefore determines neither pointwise posterior at $o=0$.
\status{ESTABLISHED}

\noindent\emph{Proof.}
The two functions differ only on the Lebesgue-null line $\{0\}\times\R$, and integrating out $y$ at $o=0$
gives $\phi(0)$ in both cases. \hfill$\square$

\noindent\emph{Proof of \Cref{lem:prob-kl-chain-rule}.}
(a) Measurability of the integrand is \Cref{cor:prob-kernel-kl-measurability} with $\mathsf E=\mathsf Y_1$.
Suppose first $Q\ll P$. Then $Q_1\ll P_1$, since $P_1(A)=0$ gives $P(A\times\mathsf Y_2)=0$. Let $r=dQ/dP$ and
$r_1=dQ_1/dP_1$. For measurable $A\subseteq\mathsf Y_1$ and $B\subseteq\mathsf Y_2$,
\[
\int_AQ_{2|1}(y_1,B)r_1(y_1)P_1(dy_1)=Q(A\times B)=\int_A\int_Br(y_1,y_2)P_{2|1}(y_1,dy_2)P_1(dy_1),
\]
so for each $B$ the integrands agree $P_1$-almost everywhere; over a countable generating algebra of
$\mathsf Y_2$ the exceptional sets form one $P_1$-null set $N$, off which both sides are finite measures in $B$
that agree on the algebra and hence everywhere. Off the $Q_1$-null set $N\cup\{r_1=0\}$, $Q_{2|1}(y_1,\cdot)$
therefore has the density $r(y_1,\cdot)/r_1(y_1)$ against $P_{2|1}(y_1,\cdot)$, and $\log r$ splits $Q$-almost
surely into $\log r_1$ and the logarithm of that density. Both summands have negative parts of $Q$-integral at
most $1/e$, as in \Cref{def:prob-log-density-ratio-convention}, so expectations add in $(-\infty,+\infty]$; the
first gives $\KL(Q_1\Vert P_1)$, and Tonelli's theorem for the positive and negative parts along
$Q=Q_1\otimes Q_{2|1}$ gives the integral of the conditional divergences. Suppose next $Q\not\ll P$. If
$Q_1\not\ll P_1$, the first term on the right is $+\infty$. If $Q_1\ll P_1$ and the conditional term were
finite, then $Q_{2|1}(y_1,\cdot)\ll P_{2|1}(y_1,\cdot)$ for $Q_1$-almost every $y_1$, so every $P$-null $C$, whose
sections are $P_{2|1}(y_1,\cdot)$-null for $P_1$-, hence $Q_1$-almost every $y_1$, would be $Q$-null. Conditional
versions are unique $Q_1$-almost everywhere when $Q_1\ll P_1$, and otherwise the right side is $+\infty$ for
every version.

(b) Quantize every block by finite refining partitions generating its sigma-algebra. On the finite product
alphabet a cell of positive $Q$-mass has marginal cells of positive $Q_b$-mass, so the first term on the right
is finite; if every marginal term is finite, every such cell has positive reference mass and the logarithm
splits cellwise, and otherwise a marginal cell of positive $Q_b$-mass and zero $\rho_b$-mass makes both sides
infinite. By \eqref{eq:prob-kernel-kl-partitions} for constant kernels, each of the three divergences is the
nondecreasing limit of its quantized values, which gives \eqref{eq:prob-kl-product-reference}.

(c) By (a), every $Q$ with $Q_1=\mu$ has $\KL(Q\Vert P)\geq\KL(\mu\Vert P_1)$. When
$\KL(\mu\Vert P_1)<\infty$, equality holds exactly when $Q_{2|1}=P_{2|1}$ $\mu$-almost everywhere,
which determines $Q$. When the marginal divergence is infinite, every member of the class has divergence $+\infty$.

(e) If $Q\ll P$ with $r=dQ/dP$, then $A_\#Q(B)=\int_{A^{-1}B}r(y)P(dy)=\int_Br(A^{-1}y')(A_\#P)(dy')$ for measurable
$B$, which is \eqref{eq:cg-rn-bimeasurable}, and integrating $\log(r\circ A^{-1})$ against $A_\#Q$ gives
$\int\log r(y)Q(dy)$. If $Q(S)>0=P(S)$, then $A(S)$ is measurable because $A^{-1}$ is, and
$A_\#Q(A(S))>0=A_\#P(A(S))$, so both divergences are $+\infty$.

(d) Form the laws $Q(dy)K(y,dz)$ and $P(dy)K(y,dz)$ on $\mathsf Y\times\mathsf Z$. By (a) with $y$ as the first
coordinate their divergence is $\KL(Q\Vert P)+\int\KL(K(y,\cdot)\Vert K(y,\cdot))Q(dy)=\KL(Q\Vert P)$. By (e) for the
coordinate swap and (a) with $z$ as the first coordinate, the same divergence is $\KL(QK\Vert PK)$ plus a term
in $[0,+\infty]$. \hfill$\square$

\section{The relative-log representation and integrability witnesses}
\label{sec:app-integrability-witnesses}

\propositionheading{Relative-log representation}{prop:elbo-relative-log-representation}
Under (H1)--(H2) of \Cref{hyp:elbo-evidence-domain}, let $M_{\theta,V,o,X}$ be the slice measure
\eqref{eq:elbo-slice-measure}. If $Q\ll M_{\theta,V,o,X}$ and $r=dQ/dM_{\theta,V,o,X}$, then
\begin{equation}
\begin{gathered}
\Lelbo_V^{\mathrm{ext}}(Q;X,o)=-\int_{\mathsf Y}\log r(y)Q(dy),\\
\int_{\{r<1\}}-\log r(y)Q(dy)=\int_{\{r<1\}}-r(y)\log r(y)M_{\theta,V,o,X}(dy)\leq\frac{p_{\theta,V}(o\given X)}{e},
\end{gathered}
\label{eq:elbo-relative-log}
\end{equation}
so the integral is well defined in $(-\infty,+\infty]$ and its negative lies in $[-\infty,+\infty)$. If
$Q\not\ll M_{\theta,V,o,X}$, then $\Lelbo_V^{\mathrm{ext}}(Q;X,o)=-\infty$. \status{ESTABLISHED}

\noindent\emph{Proof.}
Since $\boldsymbol\Pi_{\theta,V,o,X}=M_{\theta,V,o,X}/p_{\theta,V}(o\given X)$, the Radon--Nikodym chain rule gives
$dQ/d\boldsymbol\Pi_{\theta,V,o,X}=p_{\theta,V}(o\given X)r$. The negative part of $\log r$ is $Q$-integrable by the
displayed bound, because $-u\log u\leq1/e$ on $[0,1]$ and $M_{\theta,V,o,X}(\mathsf Y)=p_{\theta,V}(o\given X)$.
Hence $\KL(Q\Vert\boldsymbol\Pi_{\theta,V,o,X})=\log p_{\theta,V}(o\given X)+\int\log r(y)Q(dy)$ is a legitimate extended
identity, and substitution in \eqref{eq:elbo-extended-definition} gives the first branch. The slice measure and
the posterior have the same null sets, which gives the second. \hfill$\square$

\propositionheading{Integrability witnesses}{prop:elbo-integrability-witnesses}
The integrability hypotheses of \Cref{ch:elbo} are separated by the following witnesses.
\begin{enumerate}
\item[(a)] Clauses (H1)--(H3) of \Cref{hyp:elbo-evidence-domain} do not imply (H4). On $\R$ with a standard
Gaussian posterior and a standard Cauchy recognition law, (H1)--(H3) hold,
$\E_Q\log q=-\log4\pi$ is finite, $\E_Q\lvert\log p_{\theta,V}(o,Y\given X)\rvert=+\infty$, the divergence is
$+\infty$, and the extended bound is $-\infty$.
\item[(b)] Finite relative entropy does not imply (H4). On $\{2,3,\ldots\}$ with a trivial record, let the
posterior and the recognition law both equal $P(\{n\})=\mathrm c/(n(\log n)^2)$, with $\mathrm c$ normalizing.
The divergence is $0$, while $\sum_nP(\{n\})\lvert\log P(\{n\})\rvert=+\infty$, so both expectations in the
classical split equal $-\infty$ and their difference is undefined.
\item[(c)] Factorwise log-integrability (\cref{hyp:elbo-factorwise-integrability}) does not imply (H4). With
one latent $k\in\R$, generative factor $\mathcal N(0,1)$, and a record whose log density is bounded for
$k\in[0,1]$, let the recognition law spread mass $w_n=2^{-n}$ uniformly over disjoint intervals in $[0,1]$ of
widths $w_ne^{-1/w_n}$, whose total is below $e^{-2}$. Every factor logarithm is bounded on the support, while
$\log q=1/w_n$ on the $n$-th interval and $\E_Q\log q=\sum_nw_n/w_n=+\infty$.
\item[(d)] Factorwise log-integrability is strictly stronger than the generative half of (H4).
Take a root model coordinate $m\in\{2,3,\ldots\}$ with $p(m)=\mathrm c/(m(\log m)^2)$, a bridge
$k\given m\sim\mathcal N(0,\sigma_m^2)$ with $\sigma_m^2=p(m)^2/(2\pi)$, a trivial record, and the recognition
law equal to this law. Then $\log p(m)+\log p(k\given m)=-k^2/(2\sigma_m^2)$ has expectation $-\tfrac12$ and is
also the recognition log density, so (H4) holds, while $\E\lvert\log p(m)\rvert=+\infty$.
\end{enumerate}
\status{ESTABLISHED}

\noindent\emph{Proof.}
(a) Both laws have positive Lebesgue densities, so (H1)--(H3) hold. The log joint differs from $-y^2/2$ by a
constant and the Cauchy law has no second moment; the Cauchy differential entropy is $\log4\pi$; hence the
divergence $\E_Q\log q-\E_Q\log p$ is $+\infty$. (b) For large $n$, $P(\{n\})\leq1/n$, so
$\lvert\log P(\{n\})\rvert\geq\log n$ and the sum dominates $\mathrm c\sum_n1/(n\log n)=+\infty$. (c) The widths sum
to $\sum_n2^{-n}e^{-2^n}<e^{-2}$, so the intervals fit in $[0,1]$; the density on the $n$-th interval is
$w_n/(w_ne^{-1/w_n})=e^{1/w_n}$. (d) The bridge density at $k$ is
$(2\pi\sigma_m^2)^{-1/2}e^{-k^2/(2\sigma_m^2)}=p(m)^{-1}e^{-k^2/(2\sigma_m^2)}$, which gives the identity and
$\E\lvert\log p(m)+\log p(k\given m)\rvert=\tfrac12$; and $\E\lvert\log p(m)\rvert=+\infty$ as in (b). \hfill$\square$
